% Part: incompleteness
% Chapter: incompleteness-provability
% Section: lob-thm

\documentclass[../../../include/open-logic-section]{subfiles}

\begin{document}

\olfileid[pt]{inc}{inp}{lob}

\olsection{O teorema de L\"ob}

A sentença de G\"odel para uma teoria~$\Th{T}$ é um ponto fixo de $\lnot
\OProv[\Th{T}](y)$, O.s.\@, !!a{sentence}~$!G$ tal que
\[
\Th{T} \Proves \lnot \OProv[\Th{T}](\gn{!G}) \liff !G.
\]
Ela não é !!{derivable}, pois, se $\Th{T} \Proves !G$, então (a), pela
condição de !!{derivability}~(1),
$\Th{T} \Proves \OProv[\Th{T}](\gn{!G})$; e (b), de
$\Th{T} \Proves !G$ junto com
$\Th{T} \Proves \lnot \OProv[\Th{T}](\gn{!G}) \liff !G$, obtemos
$\Th{T} \Proves \lnot \OProv[\Th{T}](\gn{!G})$, de modo que $\Th{T}$
seria inconsistente. É natural perguntar agora qual é a situação de um ponto
fixo de $\OProv[\Th{T}](y)$, O.s.\@, de !!a{sentence}~$!H$ tal que
\[
\Th{T} \Proves \OProv[\Th{T}](\gn{!H}) \liff !H.
\]
Se ela fosse !!{derivable}, teríamos
$\Th{T} \Proves \OProv[\Th{T}](\gn{!H})$ pela condição~(1), mas a mesma
conclusão decorre da aplicação de modus ponens à equivalência acima.
Portanto, não concluímos que $\Th{T}$ seja inconsistente, ao menos não pelo
mesmo argumento usado no caso da sentença de G\"odel. Isso, evidentemente,
não mostra que $\Th{T}$ \emph{de fato} !!{derive}~$!H$.

Podemos avançar nessa questão se a generalizarmos um pouco. A direção da
esquerda para a direita da equivalência de ponto fixo,
$\OProv[\Th{T}](\gn{!H}) \lif !H$, é uma instância de um esquema geral
chamado \emph{princípio de reflexão}:
$\OProv[\Th{T}](\gn{!A}) \lif !A$. Ele recebe esse nome porque expressa,
em certo sentido, que $\Th{T}$ pode ``refletir'' sobre o que pode
!!{derive}; em essência, diz: ``Se $\Th{T}$ pode !!{derive}~$!A$, então
~$!A$ é verdadeira'', para qualquer~$!A$. É claro que isso só é verdadeiro
para teorias corretas, o que sugere que, em geral, as teorias não
!!{derive}rão todas as suas instâncias. Então, quais instâncias uma teoria
(forte o bastante e que satisfaça as condições de !!{derivability}) pode
!!{derive}? Certamente todas aquelas em que o próprio $!A$ é
!!{derivable}. E são exatamente essas, como mostra o próximo resultado.

\begin{thm}\ollabel{thm:lob}
Seja $\Th{T}$ !!a{axiomatizable} teoria que estende $\Th{Q}$, e suponha que
$\OProv[\Th{T}](y)$ seja uma fórmula que satisfaça as condições P1--P3 de
\olref[2in]{sec}. Se $\Th{T}$ !!{derive}s
$\OProv[\Th{T}](\gn{!A}) \lif !A$, então, de fato, $\Th{T}$ !!{derive}s
$!A$.
\end{thm}

Em outras palavras, se $\Th{T} \Proves/ !A$, então $\Th{T} \Proves/
\OProv[\Th{T}](\gn{!A}) \lif !A$. Esse resultado é conhecido como teorema
de L\"ob.

\begin{explain}
A heurística da prova do teorema de L\"ob é uma prova engenhosa de que Papai
Noel existe. (Se você não gosta dessa conclusão, pode substituí-la por
qualquer outra conclusão que prefira.) Ei-la:
\begin{enumerate}
\item Seja $X$ a sentença: ``Se $X$ é verdadeira, então Papai Noel
  existe''.
\item Suponha que $X$ seja verdadeira.
\item Então vale o que ela diz; O.s.\@, temos: se $X$ é verdadeira, então
  Papai Noel existe.
\item Como estamos supondo que $X$ é verdadeira, podemos concluir, por
  modus ponens a partir de (2) e~(3), que Papai Noel existe.
\item Conseguimos derivar (4), ``Papai Noel existe'', da hipótese~(2),
  ``$X$ é verdadeira''. Por prova condicional, mostramos: ``Se $X$ é
  verdadeira, então Papai Noel existe''.
\item Mas essa é simplesmente a sentença~$X$. Logo, mostramos que $X$ é
  verdadeira.
\item Mas então, pelo argumento (2)--(4) acima, Papai Noel existe.
\end{enumerate}
Uma formalização dessa ideia, substituindo ``é verdadeira'' por ``é
!!{derivable}'' e ``Papai Noel existe'' por~$!A$, produz a prova do teorema
de L\"ob. O artifício consiste em aplicar o lema do ponto fixo à
!!{formula}~$\OProv[\Th{T}](y) \lif !A$. O ponto fixo correspondente é a
sentença~$X$ do esboço anterior.
\end{explain}

\begin{proof}[Prova de \olref{thm:lob}]
Suponha que $!A$ seja !!a{sentence} tal que $\Th{T}$ !!{derive}s
$\OProv[\Th{T}](\gn{!A}) \lif !A$. Seja $!B(y)$ a
!!{formula}~$\OProv[\Th{T}](y) \lif !A$, e use o lema do ponto fixo para
encontrar !!a{sentence}~$!D$ tal que $\Th{T}$ !!{derive}s
$!D \liff !B(\gn{!D})$. Então cada uma das fórmulas seguintes é
!!{derivable} em $\Th{T}$:
\begin{align}
  & !D \liff (\OProv[\Th{T}](\gn{!D}) \lif !A) \ollabel{L-1}\\
  & \qquad \text{$!D$ é um ponto fixo de~$!B(y)$}\notag \\
  & !D \lif (\OProv[\Th{T}](\gn{!D}) \lif !A) \ollabel{L-2}\\
  & \qquad\text{de \olref{L-1}}\notag\\
  & \OProv[\Th{T}](\gn{!D \lif (\OProv[\Th{T}](\gn{!D}) \lif !A)}) \ollabel{L-3}\\
  & \qquad \text{de \olref{L-2}, pela condição P1}\notag \\
  & \OProv[\Th{T}](\gn{!D}) \lif \OProv[\Th{T}](\gn{\OProv[\Th{T}](\gn{!D}) \lif !A})
  \ollabel{L-4}\\
  &\qquad \text{de \olref{L-3}, usando a condição P2}\notag \\
  & \OProv[\Th{T}](\gn{!D}) \lif (\OProv[\Th{T}](\gn{\OProv[\Th{T}](\gn{!D})}) \lif \OProv[\Th{T}](\gn{!A})) \ollabel{L-5}\\
  &\qquad \text{de \olref{L-4}, usando P2 novamente} \notag\\
& \OProv[\Th{T}](\gn{!D}) \lif \OProv[\Th{T}](\gn{\OProv[\Th{T}](\gn{!D})}) \ollabel{L-6}\\
  & \qquad\text{pela condição de !!{derivability} P3} \notag\\
  & \OProv[\Th{T}](\gn{!D}) \lif \OProv[\Th{T}](\gn{!A}) \ollabel{L-7} \\
  &\qquad\text{de \olref{L-5} e \olref{L-6}}\notag\\
  & \OProv[\Th{T}](\gn{!A}) \lif !A \ollabel{L-8}\\
  &\qquad\text{pela hipótese do teorema} \notag\\
  & \OProv[\Th{T}](\gn{!D}) \lif !A \ollabel{L-9}\\
  &\qquad\text{de \olref{L-7} e \olref{L-8}}\notag\\
  & (\OProv[\Th{T}](\gn{!D}) \lif !A) \lif !D \ollabel{L-10}\\
  & \qquad \text{de \olref{L-1}}\notag \\
  & !D \ollabel{L-11}\\
  & \qquad\text{de \olref{L-9} e \olref{L-10}}\notag \\
  & \OProv[\Th{T}](\gn{!D}) \ollabel{L-12}\\
  & \qquad\text{de \olref{L-11}, pela condição~P1}\notag \\
  & !A \qquad\qquad\text{de \olref{L-8} e \olref{L-12}}\notag
\end{align}
\end{proof}

Com o teorema de L\"ob, há uma prova curta do segundo teorema da
incompletude (para teorias que tenham !!a{derivability} predicado que
satisfaça as condições P1--P3): se $\Th{T} \Proves
\OProv[\Th{T}](\gn{\lfalse}) \lif \lfalse$, então
$\Th{T} \Proves \lfalse$. Se $\Th{T}$ é consistente,
$\Th{T} \Proves/ \lfalse$. Logo,
$\Th{T} \Proves/ \OProv[\Th{T}](\gn{\lfalse}) \lif \lfalse$, O.s.\@,
$\Th{T} \Proves/ \OCon[\Th{T}]$. Também podemos aplicá-lo para mostrar que
~$!H$, o ponto fixo de $\OProv[\Th{T}](x)$, é !!{derivable}. Pois, como
\begin{align*}
  \Th{T} & \Proves \OProv[\Th{T}](\gn{!H}) \liff !H\\
  \intertext{em particular}
    \Th{T} & \Proves \OProv[\Th{T}](\gn{!H}) \lif !H
\end{align*}
segue-se, pelo teorema de L\"ob, que $\Th{T} \Proves !H$.

% Going in the other direction, for homework I
% may ask you to work through a short proof of L\"ob's theorem, using
% the second incompleteness theorem instead of the fixed-point lemma.

\begin{prob}
Seja $\Th{T}$ uma teoria axiomatizada computavelmente, e seja
$\OProv[\Th{T}]$ !!a{derivability} predicado para $\Th{T}$. Considere as
quatro afirmações seguintes:
\begin{enumerate}
\item Se $T \Proves !A$, então $T \Proves \OProv[\Th{T}](\gn{!A})$.
\item $T \Proves !A \lif \OProv[\Th{T}](\gn{!A})$.
\item Se $T \Proves \OProv[\Th{T}](\gn{!A})$, então $T \Proves !A$.
\item $T \Proves \OProv[\Th{T}](\gn{!A}) \lif !A$
\end{enumerate}
Sob quais condições cada uma dessas afirmações é verdadeira?
\end{prob}

\end{document}

